\section{Evaluation against the Jazz Harmony Treebank}
\label{sec:jht-evaluation}

In response to the need for an external evaluation of the proposed parser, we
implemented an evaluation bridge between our Lambek-calculus analysis system
and the Jazz Harmony Treebank (JHT) of Harasim et al.~\cite{harasim2020jazz}.
The purpose of this experiment is not merely to show that our system can parse
musical sequences, but to measure whether the harmonic relations recovered by
the system correspond to independently annotated expert analyses.

\subsection{Reference corpus}

The Jazz Harmony Treebank provides chord sequences from the iRealPro jazz
corpus together with expert harmonic tree annotations. Each annotated item
contains a chord sequence, a global key label, and one or more constituent-tree
analyses. In the JHT representation, internal tree nodes are labelled by
harmonic heads, while leaves correspond to chord events. This makes the corpus
particularly appropriate as an external reference point for our system: both
frameworks analyse the hierarchical organisation of jazz harmony, even though
they do so using different formal languages.

For the present evaluation we used the public JHT file \texttt{treebank.json}
from the \texttt{DCMLab/JazzHarmonyTreebank} repository. The evaluation script
is implemented as \texttt{jht\_eval.py}. It reads the JHT annotations, converts
JHT chord spellings to the notation accepted by our parser, runs our automatic
Lambek analysis on the same chord sequences, and writes a CSV file containing
piece-level comparison metrics. The command used for the complete run was:

\begin{verbatim}
python3 jht_eval.py \
  --jht-json /tmp/jht_treebank.json \
  --out jht_evaluation/jht_cadence_eval_full.csv
\end{verbatim}

\subsection{Notation normalisation}

Before comparison, JHT chord labels are normalised into the chord notation used
by our parser. The following deterministic substitutions are applied:

\begin{center}
\begin{tabular}{lll}
\hline
JHT notation & Normalised notation & Meaning \\
\hline
\texttt{\^{}} & \texttt{maj} & major quality, e.g. \texttt{C\^{}7} $\mapsto$ \texttt{Cmaj7} \\
\texttt{\%} & \texttt{m7b5} & half-diminished quality \\
\texttt{-} & \texttt{b} & flat accidental, e.g. \texttt{B-} $\mapsto$ \texttt{Bb} \\
\texttt{o7} & \texttt{dim7} & diminished seventh quality \\
\texttt{o} & \texttt{dim} & diminished quality \\
\hline
\end{tabular}
\end{center}

This normalisation is deliberately conservative. It does not attempt to
reinterpret the harmonic analysis supplied by JHT; it only ensures that both
systems operate over the same surface chord symbols.

\subsection{Iterative construction of the evaluation bridge}

The evaluation bridge was developed in several stages. The first version only
compared the parser's initial inferred key with the global key label supplied
by JHT. This produced a relatively low key-agreement score, because our parser
uses local tonal contexts during the derivation, while JHT provides a single
piece-level key label. Manual inspection showed that this direct comparison was
too strict: in many cases the parser correctly identified local cadential
relations, but the first or last local key did not correspond to the global
reference key of the tune.

We therefore added progressively richer key estimates. First, we extracted the
final key reached by the proof and the key implied by the last recognised
cadential relation. Then we computed the modal cadential key, namely the key
receiving the largest number of recognised cadential resolutions. This made it
possible to distinguish three different notions that are often conflated in
lead-sheet analysis: the opening local key, the most frequent local cadential
target, and the form-level key of the tune.

The first combined global-reference estimator reached 70.3\% agreement with
JHT. A manual error analysis of the remaining mismatches revealed two main
problems. First, repeated local cadences could dominate the score even when
they were part of a temporary tonicisation. Second, a spelling bug in the
dominant-implied key labels caused flat keys such as Bb and Ab to be formatted
incorrectly in some contexts. After capping repeated cadence counts and fixing
the flat-accidental formatting, agreement rose to 85.8\%.

The remaining errors were then analysed one by one. Many were caused by final
turnarounds: the last ii--V or dominant region often pointed to the next
chorus rather than to the form-level key. We therefore added a final-turnaround
detector and discounted the final cadence when it behaved as a turnaround.
The estimator was also extended to report a ranked list of candidate keys and
a confidence label. This made the residual errors inspectable rather than
opaque: the system no longer simply returned one key, but also reported how
close the competing tonal centres were.

Finally, we observed a recurrent pattern in the remaining mismatches: several
standards begin on the subdominant. In such cases the opening chord can mislead
a purely opening-based key prior. We therefore added a conditional
formal-cadence prior. The system searches for the last complete authentic
cadence resolving onto a stable, non-dominant chord. If the opening key is IV
of that formal cadence, the opening prior is reduced and the formal cadence is
given additional weight. This correction recovered cases such as
\emph{Just Friends}, whose chart opens in C but whose form-level key is G, and
\emph{I Can't Believe That You're in Love with Me}, whose chart opens in F but
whose form-level key is C. The rule is deliberately conditional: an earlier,
broader version over-weighted stable chords inside turnarounds and introduced
regressions, whereas the final version only applies when the opening-IV
relation is detected.

\subsection{Shared cadence vocabulary}

A direct tree-to-tree comparison would be misleading at this stage because the
two systems produce different kinds of structures. The JHT represents harmonic
constituency trees, whereas our system produces Lambek proof trees whose depth
also reflects derivational steps and modal-key transitions. Therefore, we first
translate both analyses into a shared intermediate vocabulary of cadential and
functional events.

The shared vocabulary currently contains the following event classes:

\begin{center}
\begin{tabular}{ll}
\hline
Event class & Description \\
\hline
\texttt{perfect\_authentic} & dominant-function resolution by descending fifth, e.g. V--I \\
\texttt{plagal} & subdominant-to-tonic relation, e.g. IV--I \\
\texttt{ii\_to\_V} & predominant-to-dominant relation, e.g. ii--V \\
\texttt{tritone\_sub} & tritone-substitute dominant resolution, e.g. bII7--I \\
\texttt{backdoor} & dominant seventh relation by whole-step descent into the target \\
\texttt{deceptive} & dominant motion into a minor target by the same root relation \\
\texttt{descending\_fifth} & generic descending-fifth root motion not otherwise classified \\
\hline
\end{tabular}
\end{center}

For our system, events are extracted from adjacent chord transitions after
normalisation, using the same dominant-quality, minor-seventh, diminished, and
tritone-substitution heuristics as the parser. For JHT, events are extracted
from the annotated constituent tree: each internal node is treated as a
harmonic head, and each non-identical child label is interpreted as a dependent
standing in a relation to that head. This maps both analyses into the same
cadence ontology while preserving the fact that JHT relations are hierarchical
and our extracted events are obtained from the parser's chord-level analysis.

\subsection{Metrics}

The evaluation reports three families of metrics.

First, we measure surface alignment:
\[
\text{leaf-count match} =
\begin{cases}
1 & \text{if the number of JHT leaves equals the number of input chords},\\
0 & \text{otherwise}.
\end{cases}
\]
We also compute leaf sequence accuracy over the comparable prefix of the two
sequences after chord normalisation. Since JHT trees may introduce headed or
repeated nodes that do not correspond one-to-one to surface chord events, we
also report a longest-common-subsequence (LCS) alignment score between the JHT
leaf sequence and the input chord sequence. This gives a less brittle measure
of whether the two systems are analysing the same harmonic material.

Second, we measure key agreement. The global key supplied by JHT is compared
with several key estimates produced by our system: the initial inferred key,
the final key label reached by the proof, the key implied by the last complete
cadential relation in the sequence, and the modal cadential key, defined as the
key receiving the largest number of recognised cadential resolutions. Finally,
we compute a global reference key estimate. This estimate combines the initial
key, final key, last-cadence key, the number of cadential resolutions per key,
and a chart-level prior inferred from the opening non-dominant chord. Cadential
counts are capped so that repeated local tonicisations cannot linearly
overwhelm the form-level evidence. A final tie-break preserves the mode of the
opening tonic when the competing estimates have the same pitch-class root.
The estimator also detects likely final turnarounds, discounts the final
cadence in those cases, and reports both a ranked list of candidate keys and a
confidence label. In addition to exact key agreement, we also report a
related-key agreement score. This score treats parallel keys, relative
major/minor pairs, and the common minor-to-flat-VI relation (for example
C minor and Ab major) as closely related tonal descriptions rather than full
errors. This separates the local key used by the Lambek derivation from a more
form-level key estimate suitable for comparison with JHT.

Third, we measure cadence agreement. Let $P$ be the multiset of cadence events
predicted by our system and $R$ the multiset of cadence events extracted from
the JHT tree. The overlap is computed category-wise:
\[
\mathrm{overlap}(P,R) = \sum_{c} \min(P_c, R_c).
\]
We then compute:
\[
\mathrm{precision} = \frac{\mathrm{overlap}(P,R)}{|P|},
\qquad
\mathrm{recall} = \frac{\mathrm{overlap}(P,R)}{|R|},
\qquad
F_1 = \frac{2PR}{P+R}.
\]
These scores are reported per annotated tree and also aggregated over the full
evaluation set.

\subsection{Results}

The complete JHT file contains 1170 pieces. Of these, 155 annotated trees were
directly comparable by the current evaluation script. No pieces were skipped
because of parser errors.

\begin{center}
\begin{tabular}{lr}
\hline
Quantity & Value \\
\hline
JHT pieces inspected & 1170 \\
Comparable annotated trees & 155 \\
Skipped pieces & 0 \\
Key match rate & 52.3\% \\
Final-key match rate & 52.3\% \\
Last-cadence key match rate & 68.4\% \\
Modal-cadential key match rate & 65.2\% \\
Representative-cadential key match rate & 72.9\% \\
Global-reference key match rate & 94.8\% \\
Global-reference related-key match rate & 96.8\% \\
Leaf-count match rate & 43.2\% \\
Mean leaf sequence accuracy & 99.9\% \\
Mean leaf LCS $F_1$ & 98.0\% \\
Median leaf LCS $F_1$ & 98.4\% \\
Mean cadence precision & 88.3\% \\
Mean cadence recall & 84.6\% \\
Mean cadence $F_1$ & 85.6\% \\
Median cadence $F_1$ & 87.0\% \\
Micro cadence precision & 90.0\% \\
Micro cadence recall & 84.9\% \\
Micro cadence $F_1$ & 87.4\% \\
\hline
\end{tabular}
\end{center}

The high leaf sequence accuracy and leaf LCS $F_1$ indicate that the
normalisation layer aligns the chord vocabulary of the two systems reliably.
The lower exact leaf-count match rate should therefore not be interpreted as a
failure of chord alignment. It is a strict structural diagnostic: several JHT
trees contain inserted, repeated, or headed nodes that do not correspond
one-to-one to surface chord tokens. This is also visible in the
structural-depth comparison. The mean depth difference is positive ($+32.7$),
confirming that our Lambek proof trees are systematically deeper than the JHT
constituent trees. This does not necessarily indicate analytical error; it
reflects the fact that our representation encodes derivational and modal
operations that are not represented as separate depth-increasing steps in JHT.

The cadence-level comparison is more directly informative. Across the 155
comparable trees, the system obtains a mean cadence $F_1$ of 85.6\% and a
micro-averaged cadence $F_1$ of 87.4\%. This suggests that, although the two
systems differ structurally, they recover substantially overlapping harmonic
relations when projected onto a common cadence vocabulary.

\begin{center}
\begin{tabular}{lrr}
\hline
Cadence class & JHT count & Our system count \\
\hline
\texttt{perfect\_authentic} & 995 & 904 \\
\texttt{plagal} & 138 & 113 \\
\texttt{ii\_to\_V} & 735 & 743 \\
\texttt{tritone\_sub} & 210 & 181 \\
\texttt{backdoor} & 55 & 29 \\
\texttt{deceptive} & 19 & 6 \\
\texttt{descending\_fifth} & 138 & 184 \\
\hline
\end{tabular}
\end{center}

The closest correspondence is observed for \texttt{ii\_to\_V}, where the
aggregate counts are nearly identical (735 in JHT and 743 in our system).
Perfect authentic cadences and plagal relations are also recovered in broadly
similar proportions. Larger discrepancies occur for deceptive and backdoor
relations, which are comparatively sparse in the corpus and depend more
strongly on analytical interpretation. The system also produces more generic
\texttt{descending\_fifth} events than JHT, suggesting that some relations that
JHT treats through a more specific hierarchical label are currently captured by
our evaluation bridge as general fifth motion.

\subsection{Residual mismatch analysis}

After applying the global-reference estimator and the related-key criterion,
only 5 of the 155 comparable annotated trees remain unmatched. These residual
cases are shown in Table~\ref{tab:jht-residual-mismatches}. Importantly, they
are not homogeneous errors. Some are low-confidence cases with sparse or
unstable evidence, while others appear to reflect chart-version or metadata
differences in the reference key label.

\begin{table}[t]
\centering
\begin{tabular}{lllll}
\hline
Tune & JHT key & Our key & Confidence & Interpretation \\
\hline
\emph{Light Blue} & C & F & low & competing F/G/C centres \\
\emph{Minor Strain} & G & d & low & sparse evidence \\
\emph{Boy Next Door} & C & Bb & medium & likely metadata/version issue \\
\emph{Remember} & Ab & Db & medium & likely chart-version issue \\
\emph{A Beautiful Friendship} & Ab & C & low & little Ab evidence in chart \\
\hline
\end{tabular}
\caption{Residual key mismatches after related-key filtering.}
\label{tab:jht-residual-mismatches}
\end{table}

The case of \emph{The Boy Next Door} is particularly illustrative. The chord
sequence analysed from JHT begins:
\[
\mathrm{Bbmaj7}\;|\;\mathrm{G7}\;|\;\mathrm{Cm7}\;|\;\mathrm{F7}\;|\;\mathrm{Bbmaj7},
\]
which strongly supports Bb major as a local and form-level centre
(\(\mathrm{Imaj7}\)--\(\mathrm{VI7}\)--\(\mathrm{ii7}\)--\(\mathrm{V7}\)--\(\mathrm{Imaj7}\)).
Our system accordingly selects Bb, whereas the JHT metadata lists C. In this
case the disagreement is therefore better interpreted as a likely metadata or
chart-version discrepancy rather than a parser failure.

The other residual cases are similarly difficult to adjudicate automatically.
For \emph{Light Blue}, \emph{Minor Strain}, and \emph{Remember}, reliable
online lead-sheet evidence is sparse or inconsistent. For \emph{A Beautiful
Friendship}, the analysed chart gives stronger evidence for C, G, F, and D
than for Ab. We therefore do not claim that all residual mismatches are JHT
errors. Rather, the residual 3.2\% disagreement should be understood as a
mixture of low-confidence tonal ambiguity, chart-version differences, and
reference metadata uncertainty.

\subsection{Interpretation}

These results provide evidence that the proposed parser is not merely producing
well-formed derivations, but is recovering a substantial portion of the
cadential structure encoded in an independent expert-annotated jazz harmony
resource. The evaluation also highlights where the system differs from JHT.
The main divergences concern global key inference, sparse cadence categories
such as deceptive and backdoor relations, and the structural depth of the
resulting analyses.

The key-agreement result should be interpreted carefully. Jazz standards often
begin away from the global tonic or contain local tonicisations at the start of
a form. Our system currently prioritises early ii--V and local cadential
evidence when selecting the initial tonal context. This behaviour is musically
motivated, but it can diverge from JHT's global key label. Comparing JHT's key
with the key implied by the last complete cadence increases agreement from
52.3\% to 68.4\%. This supports the interpretation that many mismatches are not
failures of chord parsing, but differences between local operative keys and
global form-level key labels. By contrast, the modal cadential key reaches
65.2\% agreement, while the representative-cadential key reaches 72.9\%. These
values show that cadential evidence is informative, but also that many jazz
standards contain multiple local cadential centres, so the most frequent
cadential target may still reflect an extended tonicisation rather than the
form-level key. The combined global-reference estimate obtains the strongest
result, 94.8\% agreement with JHT. This indicates that a more appropriate
comparison should not use the parser's first local key directly, but should
aggregate opening, closing, modal, turnaround, and cadential evidence over the
whole analysed sequence. Consequently, residual key mismatches do not always
imply an incorrect local analysis; rather, they show that our parser and JHT
sometimes choose different levels of tonal description. If closely related
tonalities are treated as compatible descriptions, agreement increases to
96.8\%.

\subsection{Limitations and future work}

The present evaluation is intentionally functional rather than purely
structural. It does not yet compute a full tree edit distance or labelled span
$F_1$ between JHT constituent trees and Lambek proof trees. Such a metric would
require an additional projection from Lambek derivations to constituent spans.
The current cadence ontology is therefore best understood as a common
functional layer: it asks whether both systems identify similar harmonic
relations, not whether they construct identical trees.

Future work will extend this evaluation in two directions. First, we will add a
span-based comparison by projecting each Lambek proof into a binary constituent
tree over chord positions and comparing it with JHT using labelled and
unlabelled span $F_1$. Second, we will refine the cadence ontology so that
generic descending-fifth events can be redistributed into more specific
functional categories when the local key context supports such a distinction.

Even in its present form, the evaluation addresses the main concern raised by
the reviewer: the system is now tested against an external, expert-annotated
jazz harmony resource, and the reported results quantify both agreement and
divergence at the level of musically interpretable cadential functions.
